\subsection{Work Remaining (Layer IV: Search- and Documentation-Level Poisoning of Coding Agents)}

\noindent \textbf{Position in the thesis.}
This remaining work constitutes the fourth layer of the dissertation.
Layers I and II study runtime defenses for smart contracts, including invariant enforcement and control-flow integrity.
Layer III studies harmful artifacts in code suggested by large language models (LLMs).
Layer IV moves one step further upstream by examining an external attack surface:
search- and documentation-level poisoning of coding agents.

\paragraph{Motivation.}
Modern coding agents increasingly rely on external resources such as official documentation,
API references, tutorials, and search results, when solving software engineering tasks.
While this substantially improves utility, it also introduces a new attack surface:
an adversary may influence agent behavior by manipulating the external content that the agent retrieves and trusts,
without changing model weights or injecting harmful prompts.

The core question is therefore end-to-end:
when poisoned content is present in retrieval channels, how often is it
(1) retrieved,
(2) treated as credible,
and (3) translated into concrete harmful actions?
This question is especially important because human developers do not treat all external sources equally.
Experienced developers often distinguish between mature and trustworthy libraries,
less reliable packages, unofficial tutorials, and low-quality discussions,
and they adjust their reliance accordingly. 
Experienced human developers might choose to ignore or discount 
less reliable sources, and implement their own checks or artifacts despite 
the extra effort and engineering time that doing so requires. 

However, it remains unclear whether coding agents exhibit a similar notion of source reliability,
or whether poisoned search results or documentation can alter their preferences 
and induce unsafe decisions.
This work therefore studies the interaction between an agent's prior tendencies and the external evidence it retrieves,
with a focus on when that interaction becomes security-critical.

\paragraph{Problem statement.}
This work studies \emph{external-content poisoning against coding agents}.
The attacker controls a subset of externally visible content,
such as documentation pages, API references, tutorial fragments, issue discussions, 
package descriptions, or search snippets,
while the victim is a coding agent that retrieves such content during task solving.
The attacker's goal is to induce harmful outcomes, including insecure patches,
plausible but incorrect patches, unsafe dependency choices, or unsafe recommendations in the final output.
The central research challenge is to measure and explain compromise across the full decision pipeline:
\emph{retrieval}, \emph{trust}, and \emph{action}.

\paragraph{Research questions.}
\begin{itemize}
    \item \textbf{RQ1.}
    To what extent are open-source coding agents (e.g., OpenCode) vulnerable to poisoned search results
    and documentation pages during realistic software-engineering tasks that require external lookup?

    \item \textbf{RQ2.}
    How does compromise vary with attacker budget, retrieval rank, corroborating poisoned sources,
    source reputation, and detectability constraints?
    In particular, do coding agents distinguish between sources of different reliability,
    or can poisoning override such preferences?

    \item \textbf{RQ3.}
    Which practical and deployable defenses reduce compromise
    while preserving benign task performance and agent utility?
\end{itemize}

\paragraph{Threat model.}
We assume a black-box attacker, meaning an attacker who cannot directly modify the internal model or agent implementation.
More specifically, the attacker cannot modify model weights, namely the learned internal parameters of the model,
agent system prompts, namely the hidden instructions that guide the agent's behavior, or the agent's internal implementation.
However, the attacker can control part of the external content visible through search and documentation channels
by publishing or manipulating publicly visible web content.
The victim is a coding agent that is allowed to browse or retrieve such content while solving a task.
The primary outcomes of interest are:
(1) security-degrading patches,
(2) unsafe dependency selections or API usage,
and
(3) unsafe recommendations or citations in final outputs.

\paragraph{Experimental design.}
The empirical study will evaluate at least two open-source coding-agent stacks
(e.g., an OpenCode-style system and a SWE-agent-style system)
on task suites that genuinely require external lookup.
To ensure reproducibility, the main benchmark will be executed over a controlled local retrieval corpus
that combines benign resources with attacker-controlled content, where a retrieval corpus is the document collection from which the agent searches and retrieves information.
This design enables systematic manipulation of rank, visibility, source agreement, and source reputation.

The evaluation will consider several poisoning strategies, including:
(1) \emph{rank poisoning}, where attacker-controlled content is made highly visible;
(2) \emph{content poisoning}, where retrieved documents contain misleading or adversarial technical guidance; and
(3) \emph{consensus poisoning}, where multiple poisoned sources appear to corroborate one another.
Measurements will be collected at each stage of the pipeline,
including whether poisoned content is retrieved, whether it is cited or followed,
and whether it leads to harmful code actions.
The study will also measure the tradeoff between attack resistance and normal task success.

\paragraph{Defenses and expected deliverables.}
The defense study will compare lightweight and deployable mechanisms, including:
source-level filtering,
instruction-level separation,
cross-source agreement gating,
and post-generation security checking.
The expected deliverables include:
\begin{itemize}
    \item a benchmark for search- and documentation-level poisoning of coding agents,
    \item deployment-oriented guidance for reducing upstream risk in AI-assisted software development.
\end{itemize}

\paragraph{Scope control.}
To keep this dissertation layer focused and publishable,
the work targets coding agents and documentation/search poisoning only.
It does not attempt to cover all browser-agent environments,
all prompt-injection settings,
or all forms of model poisoning.
This scope aligns Layer IV with the dissertation's overall goal of developing practical 
techniques for hardening the software-development pipeline against security threats.
